
\section{Procedencia, transporte y logística}

Sea \(\bm\ell_i\) el vector petrográfico o geoquímico de un elemento y \(q_j\) una fuente candidata.

\begin{equation}
\boxed{
P(q_j\mid\bm\ell_i)
\propto
P(\bm\ell_i\mid q_j)P(q_j).
}
\end{equation}

Un prior de transporte puede expresarse mediante

\begin{equation}
\gamma^\star
=
\arg\min_{\gamma\in\Pi(\mu_Q,\mu_P)}
\int c(q,p,t)\,d\gamma(q,p),
\end{equation}

con restricciones de capacidad y cronología. No se asume optimización histórica global; se utiliza como modelo nulo.

Para una red logística \(\mathcal N(t)=(V_N,E_N)\),

\begin{equation}
0\le f_e(t)\le c_e(t)
\end{equation}

y

\begin{equation}
\frac{dS_v}{dt}
=
\sum_{e\in\mathrm{in}(v)}f_e
-
\sum_{e\in\mathrm{out}(v)}f_e.
\end{equation}

\begin{principle}[Jerarquía de viabilidad]
\begin{equation}
\boxed{
\text{posible}
\rightarrow
\text{factible}
\rightarrow
\text{escalable}
\rightarrow
\text{históricamente compatible}.
}
\end{equation}
\end{principle}

\section{Inferencia bayesiana e identificabilidad}

Sea

\begin{equation}
Z=
(X_T,P,T_b^+,T_b^-,Q,R,W,D,C,\mathcal T,\ldots)
\end{equation}

el estado latente.

Entonces

\begin{equation}
\boxed{
P(Z\mid Y,\mathcal L)
\propto
P(Y\mid Z,\mathcal L)P(Z).
}
\end{equation}

Cuando una independencia condicional aproximada sea defendible,

\begin{equation}
P(Y\mid Z)
=
\prod_m
P(Y_m\mid Z),
\end{equation}

donde las modalidades pueden incluir geometría, muografía, litología, contactos, daño, hidrología y textos \cite{kaipio2005,tarantola2005,gelman2013}.

\begin{definition}[No identificabilidad observacional]
Dos historias \(H_1,H_2\) son no identificables bajo \(Y\) si
\begin{equation}
P(Y\mid H_1)
\approx
P(Y\mid H_2).
\end{equation}
\end{definition}

\begin{proposition}
Si dos historias poseen igual prior y generan exactamente la misma distribución de observaciones, sus odds posteriores permanecen iguales.
\end{proposition}

La consecuencia es

\begin{equation}
\boxed{
\text{evidencia faltante}
\not\Rightarrow
\text{fuerza desconocida}.
}
\end{equation}

Debe producir el estado

\begin{equation}
\boxed{
\text{NON-IDENTIFIABLE}.
}
\end{equation}

\section{Diseño activo de nuevas observaciones}

Para una medición futura \(m\),

\begin{equation}
\boxed{
\mathrm{EIG}(m)
=
\E_m
\left[
\KL
\left(
P(H\mid Y,m)
\|
P(H\mid Y)
\right)
\right].
}
\end{equation}

Esto permite priorizar muografía, fotogrametría, muestreo litológico o geofísica por su capacidad esperada de discriminar historias.

\section{Motores directo, inverso y contrafactual}

El sistema separa tres operadores:

\begin{align}
\forward &: H\rightarrow \widehat Y_H,\\
\inverse &: Y\rightarrow P(H\mid Y),\\
\counter &: H\rightarrow \text{predicciones falsables}.
\end{align}

Para una historia \(H\),

\begin{equation}
\boxed{
\widehat Y_H
=
\obs
\circ
\post
\circ
\build(H).
}
\end{equation}

Un funcional de discrepancia multimodal puede escribirse

\begin{equation}
J(H)
=
\sum_m
\lambda_m
d_m(Y_m,\widehat Y_{H,m}).
\end{equation}

Una historia que no genera observaciones potencialmente falsables no constituye todavía una reconstrucción científica fuerte.
